\documentclass{article}
% stolen from https://www.cs.cmu.edu/~15131/f17/topics/latex/
\usepackage{amsmath, amsfonts, amsthm, amssymb}  % Some math symbols
\usepackage{mathtools} % defines \coloneqq
\usepackage{enumerate}
\usepackage{fullpage}

\usepackage[x11names, rgb]{xcolor}
\usepackage{tikz}
\usepackage{graphicx}

\usepackage{algorithm}
\usepackage{algpseudocode}

\setlength{\parindent}{0pt}
\setlength{\parskip}{5pt plus 1pt}
\pagestyle{empty}

\def\indented#1{\list{}{}\item[]}
\let\indented=\endlist

\newcounter{questionCounter}
\newcounter{partCounter}[questionCounter]

\newenvironment{question}[2][\arabic{questionCounter}]{%
    \addtocounter{questionCounter}{1}%
    \setcounter{partCounter}{0}%
    \vspace{.25in} \hrule \vspace{0.5em}%
        \noindent{\bf #2}%
    \vspace{0.8em} \hrule \vspace{.10in}%
}{}

\renewenvironment{part}[1][\alph{partCounter}]{%
    \addtocounter{partCounter}{1}%
    \vspace{.10in}%
    \begin{indented}%
       {\bf (#1)} %
}{\end{indented}}


\providecommand{\abs}[1]{\left\vert#1\right\vert}
\providecommand{\norm}[1]{\left\Vert#1\right\Vert}
\providecommand{\pnorm}[2]{\left\Vert#1\right\Vert_{L^{#2}}}
\providecommand{\pnormspace}[3]{\left\Vert#1\right\Vert_{L^{#2}(#3)}}
\providecommand{\pqnorm}[3]{\left\Vert#1\right\Vert_{L^{#2,#3}}}
\providecommand{\pqnormspace}[4]{\left\Vert#1\right\VeFrt_{L^{#2,#3}(#4)}}
\providecommand{\lz}[4]{\left\Vert#1\right\Vert_{L^{#2,#3}\log^{#4}{L}}}
\providecommand{\lzspace}[5]{\left\Vert#1\rig\ht\Vert_{L^{#2,#3}\log^{#4}{L}(#5)}}
\providecommand{\wlt}[1]{\pqnorm{#1}{2}{\infty}}
\providecommand{\Rn}[1]{\mathbb{R}^{#1}}
\providecommand{\intrn}[1]{\int_{\Rn{n}}#1 \,dx}
\providecommand{\csubset}{\subset\subset}
\providecommand{\wnorm}[1]{\lvert \mspace{-1.8mu} \lvert
\mspace{-1.8mu} \lvert #1 \rvert \mspace{-1.8mu} \rvert
\mspace{-1.8mu} \rvert_{L^{2,\infty}}}

\providecommand{\br}[1]{\langle #1 \rangle}
\providecommand{\snorm}[2]{\left\Vert#1\right\Vert_{H^{#2}}}
\providecommand{\snormspace}[3]{\left\Vert#1\right\Vert_{H^{#2}({#3})}}
\providecommand{\sd}[1]{\mathcal{D}_{#1}}
\providecommand{\se}[1]{\mathcal{E}_{#1}}
\providecommand{\sdb}[1]{\bar{\mathcal{D}}_{#1}}
\providecommand{\seb}[1]{\bar{\mathcal{E}}_{#1}}
\providecommand{\snorm}[2]{\left\Vert#1\right\Vert_{H^{#2}}}
\providecommand{\snormspace}[3]{\left\Vert#1\right\Vert_{H^{#2}({#3})}}
\providecommand{\ns}[1]{\norm{#1}^2}
\providecommand{\pns}[2]{\norm{#1}^2_{L^{#2}}}
\providecommand{\dbm}[1]{\bar{D}_m^{#1}}
\providecommand{\ip}[2]{\left(#1,#2\right)}
\providecommand{\iph}[3]{\left(#1,#2\right)_{\mathcal{H}^{#3}}}
\providecommand{\hn}[2]{\norm{#1}_{\mathcal{H}^{#2}}}
\providecommand{\carrow}[1]{\overset{#1}\to}

\def\wstar{\overset{*}{\rightharpoonup}}
\def\nab{\nabla}
\def\dt{\partial_t}
\def\hal{\frac{1}{2}}
\def\ep{\varepsilon}
\def\vchi{\text{\large{$\chi$}}}
\def\ls{\lesssim}
\def\p{\partial}
\def\pa{\partial^\alpha}
\def\sg{\mathbb{D}}
\def\db{\bar{D}}
\def\da{\Delta_{\mathcal{A}}}
\def\naba{\nab_{\mathcal{A}}}
\def\diva{\diverge_{\mathcal{A}}}
\def\Sa{S_{\mathcal{A}}}
\def\Hsig{{_0}H^1_\sigma(\Omega)}
\def\H1{{_0}H^1(\Omega)}
\def\iid{\stackrel{\text{IID}}{\sim}}

\def\F{\mathcal{F}}

\def\a{\mathcal{A}}
\def\E{\mathbb{E}}
\def\f{\mathcal{F}_{2N}}
\def\F{\mathbb{F}}
\def\g{\mathcal{G}_{2N}}
\def\h{\mathcal{H}}
\def\i{\mathcal{I}}
\def\il{\i_{\lambda}}
\def\k{\mathcal{K}}
\def\n{\mathcal{N}}
\def\N{\mathbb{N}}
\def\P{\mathbb{P}}
\def\Q{\mathbb{Q}}
\def\R{\mathbb{R}}
\def\C{\mathbb{C}}
\def\w{\mathcal{W}}
\def\x{\mathcal{X}}
\def\y{\mathcal{Y}}
\def\z{\mathcal{Z}}
\def\Z{\mathbb{Z}}
\def\cf{\mathbb{C}(\F)}

%avg integral
\def\Xint#1{\mathchoice
{\XXint\displaystyle\textstyle{#1}}%
{\XXint\textstyle\scriptstyle{#1}}%
{\XXint\scriptstyle\scriptscriptstyle{#1}}%
{\XXint\scriptscriptstyle\scriptscriptstyle{#1}}%
\!\int}
\def\XXint#1#2#3{{\setbox0=\hbox{$#1{#2#3}{\int}$ }
\vcenter{\hbox{$#2#3$ }}\kern-.6\wd0}}
\def\intbar{\Xint-}

%upper and lower ints
\def\uint{\mathchoice%
    {\mkern13mu\overline{\vphantom{\intop}\mkern7mu}\mkern-20mu}%
    {\mkern7mu\overline{\vphantom{\intop}\mkern7mu}\mkern-14mu}%
    {\mkern7mu\overline{\vphantom{\intop}\mkern7mu}\mkern-14mu}%
    {\mkern7mu\overline{\vphantom{\intop}\mkern7mu}\mkern-14mu}%
  \int}
\def\lint{\mkern3mu\underline{\vphantom{\intop}\mkern7mu}\mkern-10mu\int}
\def\gchar{\ensuremath{\mathrel{\mathrm{char}}}}

\def\L{\mathcal{L}}

\DeclareMathOperator{\curl}{curl}
\DeclareMathOperator{\adj}{adj}
\DeclareMathOperator{\tr}{tr}
\DeclareMathOperator{\diverge}{div}
\DeclareMathOperator{\supp}{supp}
\DeclareMathOperator{\dist}{dist}
\DeclareMathOperator{\card}{card}
\DeclareMathOperator{\vol}{vol}
\DeclareMathOperator*{\argmax}{arg\,max}
\DeclareMathOperator*{\argmin}{arg\,min}


\newenvironment{soln}{\paragraph{Solution:}}{\hfill$\square$}

\theoremstyle{definition}
\newtheorem{lem}{Lemma}[section]
\newtheorem{defn}[lem]{Definition}
\newtheorem{cor}[lem]{Corollary}
\newtheorem{prop}[lem]{Proposition}
\newtheorem{thm}[lem]{Theorem}
\newtheorem{remark}[lem]{Remark}
\newtheorem{example}[lem]{Example}

\newtheorem*{prop*}{Proposition}
\newtheorem*{thm*}{Theorem}
\newtheorem*{def*}{Definition}
\newtheorem*{lem*}{Lemma}
\usepackage{pdfpages}

\title{S\&DS242 HW7}
\author{Zhangir Azerbayev}
\date{Spring 2022}

\begin{document}
\maketitle

\begin{question}{Problem 1}
    \begin{part}
        We have that 
        \[\mu = \int_{\theta}^\infty xe^{x-\theta}dx = 1+\theta\]
        so we get the method of moments estimate by setting $\bar{X}=1+\theta$, which gives us $\theta=\bar{X}-1$. 
    \end{part}
    \begin{part}
        Notice that if $\theta<\min X_i$, we get 0 likelihood. Therefore we can constrain ourselves to looking for $\argmax_\theta L(\theta)$ where $\theta\leq \min X_i$. Applying this constraint, we get that
        \begin{align*}
            \argmax_\theta L(\theta) &=\argmax_\theta l(\theta)\\
                                     &= \argmax_\theta\sum_i \log e^{-(X_i-\theta)}\\
                                     &= \argmax_\theta\sum_i \theta - X_i\\
                                     &= \argmax_\theta n\theta - \sum_i X_i\\
                                     &= \argmax_\theta\theta
        \end{align*}
        Therefore $\hat{\theta}_{MLE} = \min_i X_i$. 
    \end{part}
\end{question}
\begin{question}{Problem 2}
    \begin{part} 
        Reasonable estimators are $\hat{p}=\frac{\sum_iX_i}{n}$ and $\hat{q}=\frac{\sum_iY_i}{m}$. A plug-in estimator for the log-odds ratio is 
        \[g(\hat{p}, \hat{q}) = \log\left(\frac{\frac{\hat{p}}{1-\hat{p}}}{\frac{\hat{q}}{1-\hat{q}}}\right)\]
    \end{part}
    \begin{part}
        To simplify our calculations, write
        \[g(p, q) = \log p -\log(1-p) -\log q + \log (1-q). \]
        Applying the Taylor expansion we get that 
        \[g(\hat{p}, \hat{q}) \approx +\frac{\hat{p}-p}{p(1-p)} + \frac{\hat{q}-q}{q(1-q)}\]
        By the CLT, we get that $\sqrt{n}(\hat{p}-p) \to \mathcal{N}(0, p(1-p))$ and $\sqrt{m}(\hat{q}-q)\to \mathcal{N}(0, q(1-q))$. Asymptotically, both error terms are independent Gaussians, so their sum is a Gaussian with appropriate parameters. Thus in the limit 
        \[g(\hat{p}, \hat{q}) \sim \mathcal{N}\left(g(p, q), \frac{1}{np(1-p)} + \frac{1}{mq(1-q)}\right). \]
    \end{part}
\end{question}

\begin{question}{Problem 3}
    \begin{part}
        We have that 
        \[\frac{\sqrt{n}(\hat{p}-p)}{\sqrt{\hat{p}(1-\hat{p}}}\to\mathcal{N}(0,1)\]
    \end{part}
    as $n\to\infty$. Therefore our confidence interval is
    \[\hat{p}-sqrt{\frac{\hat{p}(1-\hat{p})}{n}}z(0.025)\leq p\leq \hat{p}+\sqrt{\frac{\hat{p}(1-\hat{p})}{n}}z(0.025)\]
    \begin{part}
        We are inside the interval in question whenever 
        \[n(\hat{p}-p)^2 \leq p(1-p)z(\alpha/2)^2, \]
        or equivalently when 
        \[(n+z(\alpha/2)^2)p^2 - (2n\hat{p}+z(\alpha/2)^2)p + n\hat{p}^2\leq 0. \]
        Notice the left-hand side is a quadratic equation with positive leading coefficient, so applying the quadratic formula we get a confidence interval of 
        \[\frac{\hat{p}+\frac{z(0.025)^2}{2n}\pm z(0.025)\sqrt{\frac{\hat{p}(1-\hat{p})}{n}+\frac{z(0.025)^2}{4n^2}}}{1+\frac{z(0.025)^2}{n}}\]
    \end{part}
\end{question}
\includepdf[pages=-]{hw7.pdf}
\end{document}
